\documentclass[10pt,a4paper]{amsart}
\usepackage[margin=19mm]{geometry}
\usepackage{amsmath,amssymb,mathtools}
\usepackage[scr=rsfs,cal=euler]{mathalfa}
\usepackage{xfrac}
\usepackage[unicode,hidelinks,bookmarks=false]{hyperref}
\newcommand{\ie}{i.e.\ }
\newcommand{\cf}{cf.\ }
\newcommand{\resp}{respectively\ }
\newcommand{\Subsection}{\subsection}
\providecommand{\nobreakdash}{\nobreak\hbox{-}}
\setlength{\emergencystretch}{2em}
\multlinegap0pt
\usepackage{amssymb}
\usepackage{xcolor}
\usepackage[normalem]{ulem}
\usepackage{dsfont}
\newtheorem{lemma}{Lemma}
\def\bR{\mathbb{R}}
\def\cD{\mathcal{D}}
\def\cF{\mathcal{F}}
\def\cS{\mathcal{S}}
\def\k{\kappa}
\title{SPDEs with time-independent\\ L\'evy colored noise}
\author{Raluca M. Balan, Jinxin Wang}
\date{Source: arXiv:2604.24914v1 (CC BY 4.0)}
\begin{document}
\maketitle

\noindent\emph{Source and licence.} SPDEs with time-independent\\ L\'evy colored noise. Version arXiv:2604.24914v1 (CC BY 4.0), distributed by arXiv under the Creative Commons Attribution 4.0 International licence, http://creativecommons.org/licenses/by/4.0/, as recorded in the arXiv licence metadata for this exact version and shown on the abstract page https://arxiv.org/abs/2604.24914v1. Full licence notice: \texttt{licenses/arxiv-2604.24914-CC-BY-4.0.txt}.\par\smallskip

\noindent\textit{Editorial scope.} Counted unit: the article's lemma lemma (source0.tex lines 1266--1269). The article's complete original proof of it is delivered with it (source0.tex lines 1272--1291, one \texttt{proof} environment of 20 lines). No mathematics is written, repaired or re-derived: the statement and the proof are the article's own lines, in the article's own notation, and no retained line is substituted or re-flowed.  No further source-local context is needed: every cross-reference of every retained line resolves inside the two delivered spans.  Nothing was omitted from any retained span, so the package carries no omission record.  No retained line contains a citation, so the wrapper carries no bibliography and no bibliography entry.  No retained line contains a figure environment, an image-inclusion command or drawing code, so no image is delivered and the package declares no asset dependency.  Editorial formatting only, in addition: an AMS \texttt{amsart} preamble that restates the article's own declarations and macros used by the retained lines, namely the environments \texttt{lemma} on separate counters, together with the standard packages those lines need.  The retained environments print in this document as Lemma 1 in order of appearance, whereas the article prints the same items with the numbering of its own full document.  The complete native source of the article as distributed in the arXiv e-print for this exact version is bundled unchanged as \texttt{sources/199256/source0.tex}.
\begin{lemma}
If Hypothesis A holds, then $\int_{\bR^d}|\cF(\phi* \widetilde{G}_{t})(\xi)|^2 \mu(d\xi)
<\infty$ for any $t>0$ and $\phi \in \cD(\bR^d)$.
\end{lemma}

\begin{proof} By Theorem A.2 and Hypothesis A, $\cF G_t \in \mathcal O_M$. Hence, there exist constants $C,k>0$ (depending on $t$) such that $|\cF G_t(\xi)| \le C(1+|\xi|^2)^k$ for all $\xi\in \bR^d$.

On the other hand, by Assumption A1.(b), $|\cF \k|^2$ is tempered, so there exists $\ell>0$ such that
\[
\int_{\bR^d} \left(\frac{1}{1+|\xi|^2}\right)^{\ell}\,|\cF \k(\xi)|^2\,d\xi < \infty.
\]
Therefore,
\begin{equation*}
\begin{aligned}
& \int_{\bR^d} |\cF(\phi*\widetilde{G}_{t})(\xi)|^2\,\mu(d\xi)
 = \int_{\bR^d} |\cF\phi(\xi)|^2\,|\cF G_{t}(\xi)|^2\,\mu(d\xi) \\
& \quad \quad \quad \le C \int_{\bR^d} |\cF\phi(\xi)|^2 (1+|\xi|^2)^{2k}\,|\cF \k(\xi)|^2\,d\xi \\
%&= C \int_{\bR^d} |\cF\phi(\xi)|^2 (1+|\xi|^2)^{2k+\ell}
%\left(\frac{1}{1+|\xi|^2}\right)^{\ell}\,|\cF \k(\xi)|^2\,d\xi \\
& \quad \quad \quad \le C \sup_{\xi\in\bR^d}\Big(|\cF\phi(\xi)|^2 (1+|\xi|^2)^{2k+\ell}\Big)
\int_{\bR^d} \left(\frac{1}{1+|\xi|^2}\right)^{\ell}\,|\cF \k(\xi)|^2\,d\xi < \infty,
\end{aligned}
\end{equation*}
where for the last line we used the fact that $\cF\phi\in \cS(\bR^d)$.
\end{proof}

\end{document}
